\section*{Chapitre \ref{ch:g9:elements-universe-earth} --- Les éléments dans l'Univers et sur la Terre}

\begin{solution}{exo:g9:elements-universe-earth:1}
L'hydrogène (environ \qty{74}{\percent}) et l'hélium (environ
\qty{25}{\percent}).
\end{solution}

\begin{solution}{exo:g9:elements-universe-earth:2}
L'oxygène dans les deux cas.
\end{solution}

\begin{solution}{exo:g9:elements-universe-earth:3}
Au cœur des étoiles, et lors de leurs explosions.
\end{solution}

\begin{solution}{exo:g9:elements-universe-earth:4}
Environ \qty{3.5}{\percent}.
\end{solution}

\begin{solution}{exo:g9:elements-universe-earth:5}
$46.6 + 27.7 = \qty{74.3}{\percent}$~: environ les trois quarts.
\end{solution}

\begin{solution}{exo:g9:elements-universe-earth:6}
Ce sont des gaz très légers~: ils se sont échappés de la Terre jeune
vers l'espace (et l'hélium ne forme aucun composé qui aurait pu le
retenir dans les roches).
\end{solution}

\begin{solution}{exo:g9:elements-universe-earth:7}
$0.23 \times 70 = \qty{16.1}{kg}$.
\end{solution}

\begin{solution}{exo:g9:elements-universe-earth:8}
Les \omterm{def:g7:atoms-and-molecules:atom}{atomes} d'oxygène sont 16 fois plus lourds que les \omterm{def:g7:atoms-and-molecules:atom}{atomes}
d'hydrogène~: même moins nombreux, ils pèsent davantage.
\end{solution}

\begin{solution}{exo:g9:elements-universe-earth:9}
Environ $0.081 \times 1000 = \qty{81}{kg}$ d'aluminium.
\end{solution}

\begin{solution}{exo:g9:elements-universe-earth:10}
L'oxygène et le calcium (et, au-delà des huit principaux éléments de la
croûte, l'hydrogène, le carbone et le phosphore y sont présents en
petites quantités).
\end{solution}

\begin{solution}{exo:g9:elements-universe-earth:11}
$40 \times 1.67 \times 10^{-27} = \qty{6.68e-26}{kg}$~;
$1.06 / 6.68 \times 10^{-26} \approx 1.6 \times 10^{25}$ \omterm{def:g7:atoms-and-molecules:atom}{atomes} de
calcium.
\end{solution}

\begin{solution}{exo:g9:elements-universe-earth:12}
Les \omterm{def:g5:irreversible-changes:chemical-change}{transformations chimiques} ne créent ni ne détruisent les \omterm{def:g7:atoms-and-molecules:atom}{atomes}~:
ceux-ci passent seulement d'un composé à un autre (roche, gaz, sucre,
corps vivant). Les \omterm{def:g7:atoms-and-molecules:atom}{atomes} de carbone de la Terre servent encore et
encore.
\end{solution}

\begin{solution}{pb:g9:elements-universe-earth:1}
\textbf{1.} $0.61 \times 70 = \qty{42.7}{kg}$.

\textbf{2.} Carbone~: $0.23 \times 70 = \qty{16.1}{kg}$~; hydrogène~:
$0.10 \times 70 = \qty{7.0}{kg}$.

\textbf{3.} $61 + 23 + 10 = \qty{94}{\percent}$.

\textbf{4.} L'eau, \ce{H2O}.

\textbf{5.} $16 \times 1.67 \times 10^{-27} = \qty{2.67e-26}{kg}$.

\textbf{6.} Carbone~: $12 \times 1.67 \times 10^{-27} = \qty{2.00e-26}{kg}$~;
hydrogène~: \qty{1.67e-27}{kg}.

\textbf{7.} 16.

\textbf{8.} Un \omterm{def:g9:inside-the-atom:nucleus}{électron} est environ 1836 fois plus léger qu'un \omterm{def:g9:inside-the-atom:proton}{nucléon}~:
les \omterm{def:g9:inside-the-atom:nucleus}{électrons} font moins d'un millième de la masse.

\textbf{9.} $7.0 / 1.67 \times 10^{-27} \approx 4.2 \times 10^{27}$
\omterm{def:g7:atoms-and-molecules:atom}{atomes} d'hydrogène.

\textbf{10.} $16.1 / 2.00 \times 10^{-26} \approx 8.0 \times 10^{26}$
\omterm{def:g7:atoms-and-molecules:atom}{atomes} de carbone.

\textbf{11.} $42.7 / 2.67 \times 10^{-26} \approx 1.60 \times 10^{27}$
\omterm{def:g7:atoms-and-molecules:atom}{atomes} d'oxygène.

\textbf{12.} Notre décompte, $1.60 \times 10^{27}$, est en accord avec
l'estimation $1.61 \times 10^{27}$~: un corps de \qty{70}{kg} contient
environ $1.6 \times
10^{27}$ \omterm{def:g7:atoms-and-molecules:atom}{atomes} d'oxygène.
\end{solution}
